\section{Сравнительный анализ}

\subsection{Сравнение способов использования библиотек}

\begin{table}[H]
\centering
\begin{tabular}{|p{0.3\textwidth}|p{0.3\textwidth}|p{0.3\textwidth}|}
\hline
\textbf{Критерий} & \textbf{Статическое связывание} & \textbf{Динамическая загрузка} \\ \hline
Время связывания & Во время компиляции & Во время выполнения \\ \hline
Размер исполняемого файла & Большой (библиотека включается в программу) & Малый (библиотека загружается отдельно) \\ \hline
Гибкость & Низкая (библиотека фиксирована) & Высокая (можно менять библиотеки без перекомпиляции) \\ \hline
Производительность & Высокая (вызовы прямые) & Незначительные накладные расходы \\ \hline
Отладка & Проще & Сложнее \\ \hline
\end{tabular}
\caption{Сравнение статического и динамического использования библиотек}
\end{table}

\subsection{Преимущества и недостатки}

\begin{itemize}
    \item \textbf{Статическое связывание}
    \begin{itemize}
        \item[+] Простота компиляции и распространения
        \item[+] Нет зависимостей от внешних библиотек
        \item[-] Большой размер исполняемого файла
        \item[-] Невозможно обновить библиотеку без перекомпиляции
    \end{itemize}
    
    \item \textbf{Динамическая загрузка}
    \begin{itemize}
        \item[+] Возможность «горячей» замены библиотек
        \item[+] Экономия памяти (одна библиотека для нескольких программ)
        \item[+] Лёгкое обновление функционала
        \item[-] Сложность отладки
        \item[-] Зависимость от наличия библиотек в системе
    \end{itemize}
\end{itemize}